\newpage
\section{Acta Número 28}
\label{sec:acta28}

\noindent \textbf{Fecha de la sesión:}\par
Lunes, 08 de Junio de 2026

\vspace*{0.5cm}
\noindent \textbf{Datos de la Sesión:}
\vspace{-0.2cm}
\begin{itemize}[noitemsep]
    \item \textbf{Modalidad:} Virtual - Google Meet.
    \item \textbf{Hora de Inicio:} 18:00
    \item \textbf{Hora de Fin:} 20:00
\end{itemize}

\vspace*{0.5cm}
\noindent \textbf{Orden del Día:}
\vspace{-0.2cm}
\begin{itemize}[noitemsep]
    \item Control de Asistencia.
    \item Socialización de las nuevas peticiones de funcionalidad remitidas por la Consultora TIS.
    \item Análisis de viabilidad técnica e impacto de las nuevas funcionalidades sobre el sistema.
    \item Planificación y asignación de tareas para la implementación de los nuevos requerimientos.
    \item Firmas y Aprobación de los Presentes.
\end{itemize}

\vspace*{0.5cm}
\noindent \textbf{Socios Fundadores Presentes:}
\vspace{-0.2cm}
\begin{enumerate}[noitemsep]
    \item Pardo Pozo Juan Diego (Jefe de Proyecto).
    \item Arnez Jaimes Mauricio.
    \item Quispe Flores Camila Belen.
    \item Ariñez Bautista Jim Gabriel.
    \item Medina Rodríguez Mateo Jhoustin.
    \item Molina Maldonado Tomas Manuel.
    \item Gutierrez Guzmán Angheline.
\end{enumerate}

\newpage
\subsection{Desarrollo del Acta}

\noindent \textbf{Control de Asistencia del Team}\par
Se procedió a la toma de asistencia a la hora acordada, corroborando la presencia de los 7 socios fundadores de la grupo-empresa Byte Busters S.R.L. Al contar con el quórum y la atención de todos los involucrados, se dio inicio formal a la sesión, convocada de manera específica para atender los nuevos requerimientos de la Consultora TIS.

\vspace*{0.5cm}
\noindent \textbf{Socialización de las Nuevas Peticiones de la Consultora TIS}\par
El Jefe de Proyecto presentó al equipo el conjunto de peticiones de nuevas funcionalidades remitidas formalmente por la Consultora TIS, surgidas tras la revisión del producto entregado en el último incremento. Se expusieron en detalle los alcances solicitados, las expectativas de comportamiento del sistema y las justificaciones funcionales que respaldan cada requerimiento. El equipo tomó nota de cada petición, asegurándose de comprender con precisión la intención y el valor que la consultora espera obtener de cada nueva característica antes de avanzar a su valoración técnica.

\vspace*{0.5cm}
\noindent \textbf{Análisis de Viabilidad Técnica e Impacto}\par
A continuación, se realizó un análisis técnico pormenorizado de cada funcionalidad solicitada, evaluando su impacto sobre la arquitectura actual del sistema. Se identificaron los módulos de \textit{Backend}, \textit{Frontend} y \textit{Base de Datos} que requerirían modificaciones o nuevas adiciones, así como las posibles afectaciones sobre los entornos desplegados en \textbf{Supabase} y \textbf{Vercel}. El equipo determinó que las peticiones eran técnicamente viables y acordó las estrategias de implementación que minimizaran el riesgo de regresiones sobre las funcionalidades ya estabilizadas y validadas.

\vspace*{0.5cm}
\noindent \textbf{Planificación y Asignación de Tareas}\par
Con la viabilidad confirmada, se procedió a desglosar las nuevas funcionalidades en tareas técnicas concretas, estimando el esfuerzo requerido y distribuyéndolas entre los integrantes según sus áreas de especialidad. Se definió un cronograma de trabajo acotado para atender estas peticiones de manera ágil, estableciendo los criterios de aceptación correspondientes y reafirmando la Definición de Hecho (DoD) para garantizar que las nuevas características alcancen los estándares de calidad exigidos por la Consultora TIS.

\vspace*{0.5cm}
\noindent \textbf{Aprobación Oficial}\par
Tras socializar las nuevas peticiones de la Consultora TIS, confirmar su viabilidad técnica y distribuir las tareas de implementación con un cronograma consensuado, el equipo manifiesta su total conformidad con el plan de trabajo establecido. Se da por aprobada la sesión y se procede con la firma de los presentes.

\newpage
\noindent \textbf{Firmas y Aprobación de los Presentes}
\begin{center}
    \noindent
    \setlength{\tabcolsep}{0pt}
    \begin{tabular}{ccc}
        \espacioFirma{assets/img/firmas/mauricio.jpg}{Mauricio Jaimes Arnez}{13751221} &
        \espacioFirma{assets/img/firmas/camila.jpg}{Camila Belen Quispe Flores}{13097244} &
        \espacioFirma{assets/img/firmas/jim.jpg}{Jim Gabriel Ariñez Bautista}{9517642} \\[1.0cm]
        \espacioFirma{assets/img/firmas/mateo.jpg}{Mateo Medina Rodríguez}{9453809} &
        \espacioFirma{assets/img/firmas/tomas.jpg}{Tomas Molina Maldonado}{8051701} &
        \espacioFirma{assets/img/firmas/angheline.jpg}{Angheline Gutierrez Guzmán}{13751815} \\[1.0cm]
    \end{tabular}
    \vspace{0.2cm}
    \begin{tabular}{cc}
        \espacioFirma{assets/img/firmas/diego.jpg}{Juan Diego Pardo Pozo}{12747783}
    \end{tabular}
\end{center}

\vspace{0.5cm}
\noindent \textbf{Fecha de última revisión:} Lunes, 08 de Junio de 2026
